\section{Discussion}
\label{sec:discussion}

\subsection{Hypothesis Testing}

The experimental design is motivated by four hypotheses.

\textbf{H1 (Group~B $>$ Group~A): Not supported.}
Group~B's mean improvement ($-0.006 \pm 0.006$) is numerically worse than
Group~A's ($-0.009 \pm 0.002$), though the difference is not meaningful given
$n=3$.
The outer loop (Level~1.5) increases search diversity---Group~B explores more
parameters than Group~A---but this diversity does not translate into larger
improvements within the 30-iteration budget.
Group~B's R1 achieved essentially zero improvement ($-0.000$), the worst
outcome of any repeat in any group.
The outer loop correctly froze stalled parameters but, having done so, the LLM
found nothing better in the remaining search space.

\textbf{H2 (Group~C $>$ Group~B): Supported.}
Group~C's mean absolute improvement ($-0.045 \pm 0.030$) is 7.5$\times$
Group~B's ($-0.006 \pm 0.006$).
Despite high variance ($\pm 0.030$), two of three repeats produced dramatic
improvements ($-0.065$, $-0.058$), and the separation between Groups~C and~B
is large relative to the within-group variance, providing meaningful evidence
that Level~2 adds value beyond Level~1.5.

\textbf{H3 (Level~2 discovers novel mechanisms autonomously): Supported.}
Across three independent repeats, Level~2 generated mechanisms from three
distinct active domains (combinatorial optimization, online learning, DOE)
without being told which domains to consider (a fourth domain, Bayesian
optimization, was attempted but reverted due to a missing dependency).
Code generation succeeded on the first attempt in all six sessions (zero
retries).
Five of six mechanisms passed import validation and were activated.

\textbf{H4 (Group~D $\approx$ Group~C: Level~1.5 is not essential when Level~2 is present): Supported with caveats.}
Group~D's mean ($-0.034 \pm 0.031$) is lower than Group~C's ($-0.045 \pm 0.030$),
but the difference is within the variance of both groups given $n=3$.
Level~2 alone is sufficient to produce large improvements in two of three
repeats; Level~1.5 does not appear to be a necessary condition.
The caveat is that D's R1 produced essentially no improvement, whereas
Group~C had no zero-improvement repeat, suggesting Level~1.5 may provide
modest robustness by enriching the search trace that Level~2 reads.

\subsection{Why Group~C's R2 Underperformed}

Group~C's R2 achieved only $-0.011$ improvement, comparable to Groups~A and~B,
despite receiving Level~2 mechanisms.
The most likely explanation is mechanism quality: R2's Level~2 generated the
Multi-Scale Bandit Proposer, which---while valid and correctly injected---may
be less effective than R1's Tabu Search Manager or R3's Orthogonal Exploration
at forcing exploration of the batch size dimension.
A secondary factor is overhead: each Level~2 research session requires
approximately 3~minutes of wall time (four LLM calls), reducing effective inner
iterations.
With two sessions per repeat, Group~C has roughly 6~minutes less search time
than Groups~A and~B, a minor but non-zero cost.

\subsection{Limitations}

\paragraph{Small sample size.}
Three repeats per group is insufficient for rigorous statistical comparison.
Group~C's standard deviation ($\pm 0.030$) is 67\% of its absolute mean, indicating
high variability.
Reliable estimates would require $n \geq 10$ repeats per group.

\paragraph{Baseline variance.}
Baseline \valbpb{} varies across repeats (1.094--1.114) due to training
randomness from data ordering and weight initialization.
Using $\Delta = \text{best} - \text{baseline}$ mitigates this but does not
eliminate it; a lower baseline gives less headroom for improvement.
Future work should use fixed random seeds or report baseline-normalized metrics.

\paragraph{Single benchmark.}
All results are on one task: GPT pretraining at 50M parameters with a 300-second
budget on RTX~5090.
Generalization to other model sizes, training budgets, or tasks is unproven.

\paragraph{Dynamic load fragility.}
A preliminary run was invalidated because the Level~2 dynamic loading pipeline
contained a \texttt{sys.modules} registration bug, causing all mechanism
injections to silently fall back to the original runner.
We fixed the bug before conducting the three reported repeats.
This episode highlights the fragility of runtime code injection: silent
fallback without error is a dangerous failure mode.

\paragraph{External dependency exposure.}
Level~2 has no constraint preventing it from importing external libraries.
One of six generated mechanisms (GP Regressor) required \texttt{sklearn}, which
was not installed.
The validate-and-revert mechanism handled this correctly, but the exposure to
arbitrary dependencies remains a reliability risk.

\paragraph{Prompt-induced domain bias.}
The Level~2 prompt explicitly suggests candidate domains (combinatorial
optimization, reinforcement learning, evolutionary algorithms, Bayesian
optimization).
This guidance is a double-edged sword: it prevents the agent from generating
irrelevant or degenerate mechanisms, but it also constrains the search space
of discoverable mechanisms to domains the prompt author anticipated.
Whether Level~2 would discover equally effective---or entirely
different---mechanisms under an unconstrained prompt remains untested.

\subsection{Future Work}

Key directions include:
(1) scaling to $n \geq 10$ repeats with fixed random seeds for statistical power;
(2) evaluating on multiple benchmarks (different model sizes, tasks, budgets)
    to assess generalization; and
(3) investigating whether Level~2's code generation quality improves with
    a richer interface specification or a test harness.
